\subsection*{一步一步建立直观}

\paragraph{1. 把一个初值问题看成待求的整条信号}

选定初始状态后，局部速度规则产生整条未来轨迹。拉普拉斯方法将这条未知信号整体送到参数域，在那里把微分操作转成乘法，同时留下初始值项。右图的初值拖动可以提醒我们：初始条件改变，解就改变，变换计算不能把它丢掉。



\paragraph{2. 为什么导数变换带着边界项}

对加权积分做分部积分，导数从未知信号转移到指数核上，因而出现乘以 $s$ 的项。下端点则贡献初始值，上端点需借助衰减条件消失。因此公式不是单纯的“把求导替换成乘法”；边界与收敛假设同样参与计算。



\paragraph{3. 代数解仍要回到原问题核对}

改成增长轨迹，观察解的长期行为怎样变化。在参数域求出的分式必须在合适的收敛域内解释，再还原时间函数，最后检查微分方程与初值。右图不是拉普拉斯计算器，而是用于核对所得解的轨迹、初值与稳定性是否合理的模型。


